\section{Problem Statement}
\label{sec:problem}

That four-part gap breaks down into five concrete limitations once we look
closely at what today's approaches actually miss:

\begin{enumerate}
  \item \textbf{Coverage gap}: no study has fully labelled the active
        \textit{Defects4J} set; existing work covers only tens to a few hundred
        bugs \citep{xuan2017better,vanderspuy2025anatomy,sobreira2018dissection}.
  \item \textbf{Structural-only labels}: control-/data-flow analyses capture
        how bugs manifest syntactically, not the defect type they belong to
        \citep{vanderspuy2025anatomy}.
  \item \textbf{Patch-focused taxonomies}: repair-action classification
        describes fix patterns, not the underlying defect mechanism
        \citep{sobreira2018dissection}.
  \item \textbf{Post-fix dependency}: ML approaches require labelled training
        data or access to post-fix code, limiting applicability to pre-triage
        scenarios \citep{thung2012automatic}.
  \item \textbf{Evaluation scarcity}: without ODC ground truth, there is no
        practical evaluation framework for large-scale automated classification
        \citep{andrade2025empirical}.
\end{enumerate}

These limitations motivate an overarching question: \emph{how can we classify
\textit{Defects4J} bugs into ODC Design/Code defect types from pre-fix evidence
alone, while also being able to evaluate that classification without
ground-truth labels?} The five research questions of
Section~\ref{sec:rqs} decompose this question into measurable parts. Our
guiding intuition is that pre-fix classification is a prediction made from
symptoms alone, while post-fix classification is a reference made with the fix
visible; agreement between the two indicates the prediction is stable enough to
justify operational use of the pre-fix signal.

%% =============================================================================
